% META: {"id":"1SPE-DERGLOBAL-CO-039","chapitre":"1SPE-DERIVATION-GLOBAL","type_objet":"corrige","exercice_id":"1SPE-DERGLOBAL-EX-039","status":"generated"}
% BEGIN-VERIFY
% from sympy import symbols, diff, solve, Rational, simplify
% x, k = symbols('x k', real=True)
% f = x**2 - 2*k*x + k
% fp = diff(f, x)
% assert fp == 2*x - 2*k
% x0 = solve(fp, x)
% assert x0 == [k]
% f_at_x0 = f.subs(x, k)
% assert simplify(f_at_x0 - (k - k**2)) == 0
% assert f_at_x0.subs(k, 0) == 0
% assert f_at_x0.subs(k, 1) == 0
% assert f_at_x0.subs(k, Rational(1,2)) == Rational(1,4)
% END-VERIFY
\begin{corrige}{1SPE-DERGLOBAL-EX-039}
\textbf{1.} On dérive $f_k(x) = x^2 - 2kx + k$ par rapport à $x$ ($k$ est une constante) :
\[
f_k'(x) = 2x - 2k
\]
On résout $f_k'(x) = 0$ :
\[
2x - 2k = 0 \iff x = k
\]
Le point stationnaire est $x_0(k) = k$.

\textbf{2.} On étudie le signe de $f_k'(x) = 2(x - k)$ :
\begin{itemize}
\item Pour $x < k$ : $f_k'(x) < 0$ (décroissant).
\item Pour $x > k$ : $f_k'(x) > 0$ (croissant).
\end{itemize}
La dérivée change de signe de $-$ à $+$ en $x = k$ : $f_k$ admet bien un \textbf{minimum} en $x_0(k) = k$.

\textbf{3.} La valeur minimale est :
\[
m(k) = f_k(k) = k^2 - 2k \cdot k + k = k^2 - 2k^2 + k = -k^2 + k
\]

\textbf{4.} On étudie $m(k) = -k^2 + k$ en tant que fonction de $k$ :
\[
m'(k) = -2k + 1
\]
On résout $m'(k) = 0$ :
\[
-2k + 1 = 0 \iff k = \frac{1}{2}
\]

Signe de $m'(k)$ :
\begin{itemize}
\item Pour $k < \frac{1}{2}$ : $m'(k) > 0$ ($m$ croissante).
\item Pour $k > \frac{1}{2}$ : $m'(k) < 0$ ($m$ décroissante).
\end{itemize}

La valeur minimale $m(k)$ est \textbf{maximale} pour $k = \dfrac{1}{2}$ :
\[
m\!\left(\frac{1}{2}\right) = -\frac{1}{4} + \frac{1}{2} = \boxed{\frac{1}{4}}
\]

C'est pour $k = \frac{1}{2}$ que le minimum de $f_k$ est le plus grand possible, avec un minimum valant $\frac{1}{4}$.
\end{corrige}
